% Part: normal-modal-logic
% Chapter: axioms-systems
% Section: derived-rules

\documentclass[../../../include/open-logic-section]{subfiles}

\begin{document}

\olfileid[fr]{nml}{prf}{der}

\olsection{Règles dérivées}

Trouver et rédiger des !!{derivation}s est souvent laborieux et
répétitif. Par exemple, on veut fréquemment passer de $!A \lif !B$ à
$\Box !A \lif \Box !B$, c'est-à-dire appliquer la règle~\RK. Il faut
pour cela appliquer \Nec, écrire l'instance appropriée de~\Ax{K},
puis appliquer~\MP. Pour passer de $!A \lif !B$ et $!B \lif !C$ à
$!A \lif !C$, il faut écrire l'instance (longue) de la tautologie
\[
(!A \lif !B) \lif ((!B \lif !C) \lif (!A \lif !C))
\]
et appliquer \MP{} deux fois. On veut aussi souvent remplacer une
sous-!!{formula} par une !!{formula} dont on connaît l'équivalence :
\iftag{prvDiamond}{$\Diamond !A$ par $\lnot\Box\lnot !A$}{$\Box !A$
  par $\lnot\Diamond\lnot !A$}, par exemple, ou bien
$\lnot\lnot !A$ par $!A$. Au lieu de rédiger chaque !!{derivation}
en entier, il est donc plus commode d'indiquer pourquoi les étapes
intermédiaires sont !!{derivable}. Rassemblons à cette fin quelques
résultats sur la !!{derivability}.

\begin{prop}
  Si $\Log{K} \Proves !A_1$, \dots, $\Log{K} \Proves !A_n$, et si $!B$
  découle de $!A_1$, \dots, $!A_n$ par la logique propositionnelle,
  alors $\Log{K} \Proves !B$.
\end{prop}

\begin{proof}
  Si $!B$ découle de $!A_1$, \dots,~$!A_n$ par la logique
  propositionnelle, alors
  \[
  !A_1 \lif (!A_2 \lif \cdots (!A_n \lif !B)\dots)
  \]
  est une instance de tautologie. En appliquant \MP{} $n$ fois, on
  obtient une !!{derivation} de~$!B$.
\end{proof}

Nous noterons l'emploi de cette proposition par~\PL.

\begin{prop}
  Si $\Log{K} \Proves !A_1 \lif (!A_2 \lif \cdots (!A_{n-1} \lif
  !A_n)\dots)$, alors $\Log{K} \Proves \Box !A_1 \lif (\Box !A_2 \lif
  \cdots (\Box !A_{n-1} \lif \Box !A_n)\dots)$.
\end{prop}

\begin{proof}
  Par récurrence sur $n$, comme dans la démonstration de la
  \olref[nor]{prop:rk}.
\end{proof}

Nous noterons l'emploi de cette proposition par~\RK. Voyons comment
ces résultats facilitent l'établissement de la !!{derivability}.

\begin{prop}
  $\Log{K} \Proves (\Box!A \land \Box!B) \lif \Box (!A \land !B)$
\end{prop}

\begin{proof}
  \begin{derivation}
    1. & $\Log{K} \Proves !A \lif (!B \lif (!A \land !B))$ & \Taut \\
    2. & $\Log{K} \Proves \Box!A \lif (\Box !B \lif \Box(!A \land !B))$ & \RK, 1\\
    3. & $\Log{K} \Proves (\Box!A \land \Box!B) \lif \Box(!A \land !B)$ & \PL, 2
  \end{derivation}
\end{proof}

\begin{prop}\ollabel{prop:rewriting}
  Si $\Log{K} \Proves !A \liff !B$ et $\Log{K} \Proves
  \Subst{!C}{!A}{q}$, alors $\Log{K} \Proves \Subst{!C}{!B}{q}$.
\end{prop}

\begin{proof}
  Exercice.
\end{proof}

\begin{prob}
  Démontrer la \olref[nml][prf][der]{prop:rewriting} en montrant, par
  récurrence sur la complexité de~$!C$, que si
  $\Log{K} \Proves !A \liff !B$, alors
  $\Log{K} \Proves \Subst{!C}{!A}{q} \liff \Subst{!C}{!B}{q}$.
\end{prob}

Cette proposition est particulièrement utile pour exprimer $\Diamond$
au moyen de $\Box$ (ou inversement), ou pour éliminer les doubles
négations à l'intérieur d'une !!{formula}. Dans la suite, nous
signalerons les applications de la \olref{prop:rewriting} par
« $!B$ à la place de $!A$ » lorsque nous remplaçons
$!C(!A)$ par~$!C(!B)$. Autrement dit, « $!B$ à la place de $!A$ »
abrège :
  \begin{derivation}
    & $\Proves !C(!A)$ \\
    & $\Proves !A \liff !B$\\
    & $\Proves !C(!B)$ & par la \olref{prop:rewriting}
  \end{derivation}
Par exemple :

\begin{prop}
  $\Log{K} \Proves \lnot\Box p \lif \Diamond \lnot p$
\end{prop}

\begin{proof}
  \iftag{prvDiamond}{% Diamond est primitif
  \begin{derivation}
    1. & $\Log{K} \Proves \Diamond \lnot p \liff \lnot\Box\lnot\lnot p$ &
    \Dual\\
    2. & $\Log{K} \Proves \lnot \Box \lnot\lnot p \lif \Diamond \lnot p$ & \PL, 1\\
    3. & $\Log{K} \Proves \lnot\Box p \lif \Diamond \lnot p$ & $p$ à la place de $\lnot\lnot p$\\
  \end{derivation}
  }{% Diamond est défini
  \begin{derivation}
    1. & $\Log{K} \Proves \lnot \Box \lnot\lnot p \lif \Diamond \lnot p$ & \Taut\\
    2. & $\Log{K} \Proves \lnot\Box p \lif \Diamond \lnot p$ & $p$ à la place de $\lnot\lnot p$\\
  \end{derivation}
  La !!{formula} de la ligne~$1$ est une instance de
  $\lnot q \lif \lnot q$, obtenue en substituant
  $\Box\lnot\lnot p$ à~$q$. Les expressions $\Diamond\lnot p$ et
  $\lnot\Box\lnot\lnot p$ sont identiques, puisque $\Diamond$ est
  défini par~$\lnot\Box\lnot$.}
\end{proof}

Dans la !!{derivation} précédente, la dernière étape, « $p$ à la
place de $\lnot\lnot p$ », abrège
  \begin{derivation}
    & $\Log{K} \Proves \lnot \Box \lnot\lnot p \lif \Diamond \lnot
    p$\\
    & $\Log{K} \Proves \lnot\lnot p \liff p$ & \Taut\\
    & $\Log{K} \Proves \lnot\Box p \lif \Diamond \lnot p$ & par la \olref{prop:rewriting}
  \end{derivation}
Les rôles de $!C(q)$, $!A$ et $!B$ dans la \olref{prop:rewriting} sont
joués ici respectivement par $\lnot \Box q \lif \Diamond \lnot p$,
$\lnot\lnot p$ et~$p$.

Lorsqu'une !!{formula} contient une sous-!!{formula}
$\lnot\Diamond !A$, on peut la remplacer par $\Box\lnot !A$ au moyen
de la \olref{prop:rewriting}, puisque
$\Log{K} \Proves \lnot\Diamond !A \liff \Box\lnot !A$. Nous
indiquerons ce remplacement, et les remplacements analogues,
simplement par « $\Box\lnot$ à la place de $\lnot\Diamond$ ».

La proposition suivante justifie les démonstrations schématiques de
!!{derivability}. La proposition précédente, par exemple, n'établit
pas seulement $\Log{K} \Proves \lnot\Box p \lif \Diamond \lnot p$,
mais aussi $\Log{K} \Proves \lnot\Box !A \lif \Diamond\lnot !A$
pour toute !!{formula}~$!A$.

\begin{prop}
  Si $!A$ est une instance de substitution de $!B$ et si
  $\Log{K} \Proves !B$, alors $\Log{K} \Proves !A$.
\end{prop}

\begin{proof}
  Il est fastidieux, mais élémentaire, de vérifier par récurrence sur
  la longueur de la !!{derivation} de~$!B$ qu'une substitution
  appliquée à toute cette !!{derivation} donne encore une
  !!{derivation} correcte. Plus précisément, les instances de
  substitution d'instances de tautologies sont encore des instances
  de tautologies, celles des instances de~\iftag{prvDiamond}{\Dual{}
    et de~}{}\Ax{K} sont encore des instances de~\iftag{prvDiamond}{\Dual{}
    et de~}{}\Ax{K}, et les applications de \MP{} et de~\Nec{} restent
  correctes lorsque l'on substitue des !!{formula}s aux variables
  propositionnelles dans les prémisses et la conclusion.
\end{proof}

\end{document}
